\section{Introduction}\label{sec:introduction}

Let $a$ be an integer and let $p$ be a prime not dividing $a$.
The multiplicative order $\ord_p(a)$ is the least positive integer $k$
such that $a^k\equiv1\pmod p$. The residue class of $a$ is a primitive
root modulo $p$ if $\ord_p(a)=p-1$. Call $a$ \emph{admissible} if
$a\ne-1$ and $a$ is not the square of an integer. These are necessary
conditions for primitive-root behavior at infinitely many primes:
the order of $-1$ is at most $2$, and a nonzero square modulo an odd
prime has order dividing $(p-1)/2$.

Artin's primitive root conjecture, formulated in 1927, asserts that
every admissible integer is a primitive root modulo infinitely many
primes; its quantitative form predicts a positive-density asymptotic
\cite[Section~17.2]{ArtinHasse08}. We prove the following lower bound.

\begin{theorem}\label{thm:primitive-root}\label{thm:main}
For every admissible integer $a$, there are constants $c_a>0$ and
$x_a\ge2$ such that, for every real $x\ge x_a$,
\[
 \#\{p\text{ prime}:x<p<2x,\ p\nmid a,\ \ord_p(a)=p-1\}
       \ge c_a\frac{x}{(\log x)^2}.
\]
\end{theorem}

The predicted counting scale is $x/\log x$; the theorem gives
$x/(\log x)^2$ in every sufficiently large dyadic interval. In
particular, it proves the infinitude assertion for each fixed
admissible integer, including the prescribed base $2$. The constants
and the threshold may depend on the base.

The analytic part of the proof establishes a separate theorem.
A finite-order Hecke character of a number field $F$ means a continuous
finite-image character of the idele class group
$F^\times\backslash\Adele_F^\times$. Write $\mu_m$ for the group of
$m$th roots of unity.

\begin{theorem}\label{thm:hecke-zero-free}
Let $F$ be a cyclotomic number field containing $\mu_{12}$.
For every finite-order Hecke character $\eta$ of $F$, the meromorphic
continuation of $L_F(s,\eta)$ has no zeros in
\[
 \re s>1-10^{-6}.
\]
\end{theorem}

The principal character is included, with its pole at $s=1$ permitted.
The width is common to all fields and characters in the statement,
with no conductor or height cutoff. In the proof, $F$ and $\eta$ are
fixed before the auxiliary data are chosen. Constants and the points
from which estimates hold may depend on them; the numerical saving
in the final Mellin argument does not. This distinction is what
allows a uniform splitting estimate for a growing family of fields.

\subsection{Historical context and the two obstructions}

Hooley proved Artin's predicted asymptotic under the Riemann hypothesis
for the Dedekind zeta functions of the Kummer fields
$\Q(\mu_n,a^{1/n})$, with $n$ squarefree~\cite{Hooley67}.
These fields encode the obstructions to maximal order. For a prime
$q$ and $p\nmid aq$, the conditions
\[
 p\equiv1\pmod q,\qquad a^{(p-1)/q}\equiv1\pmod p
\]
say that $p$ splits completely in $\Q(\mu_q,a^{1/q})$.
They hold precisely when $q$ divides the index
$(p-1)/\ord_p(a)$. Hooley's method counts and sums such splitting
conditions. Here a restricted factorization of $p-1$ confines the
possible index divisors to ranges where they can all be removed.

Classical unconditional results allow a finite choice of bases.
Gupta and Murty constructed a set of thirteen integers from three
distinct primes, at least one of which is a primitive root modulo
infinitely many primes~\cite{GuptaMurty84}. Heath-Brown proved that
at most two prime bases fail the infinitude assertion, so at least
one of $2,3,5$ satisfies it~\cite{HeathBrown86}.
Averaging over the base gives another approach. Goldfeld and Stephens
established early average theorems~\cite{Goldfeld68,Stephens69};
Klurman, Shparlinski, and Ter\"av\"ainen obtained the expected
asymptotic for almost all integers in a range as short as
$1\le a\le\exp((\log\log x)^2)$~\cite{KlurmanShparlinskiTeravainen25}.
The exceptional set may depend on $x$, so these average results
do not by themselves settle a prescribed base.

Two difficulties must be addressed for a fixed base. One is analytic:
the splitting fields vary with the possible index divisor. A standard
simultaneous zero-free region for Hecke characters associated with
an abelian extension has logarithmic width depending on the field,
conductors, and height, and permits at most one simple real
exceptional zero~\cite[Theorem~3.1]{ThornerZaman19}.
The common width in Theorem~\ref{thm:hecke-zero-free} instead gives
a common power saving. The other difficulty is the prime construction:
we need many primes for which the possible index divisors fall in
the range controlled by that saving. We now describe the two parts
of the proof and the points where they meet.

\subsection{The Hecke zero-free argument}

Fix a cyclotomic field $F$ containing $\mu_{12}$ and a finite-order
Hecke character $\eta$. The analytic construction uses cubic theta
coefficients under sextic Kummer twists. Patterson's cubic theta
series over $\Q(\sqrt{-3})$ is the model for the Gauss coefficients
and their transformations~\cite{Patterson77,Patterson77II}; the
comparison formulas are recorded by Dunn and Radziwi\l{}\l{}
\cite[Section~5.1 and Appendix~A]{DunnRadziwill24}.
The companion manuscript \emph{The Quasi-Riemann Hypothesis}
develops a coefficient reflection and an additive probe over that
field while retaining unrestricted cube factors
\cite[Part~I, Sections~5--7]{QuasiRH}. Its zero-free theorem treats
Dirichlet $L$-functions and finite-order Hecke $L$-functions over
$\Q(\sqrt{-3})$. We construct the required mechanisms over each
fixed $F$, using the cubic theta representation of Kazhdan and
Patterson and a normalized global Whittaker factorization consistent
with their correction~\cite{KP84,KP85Erratum}.

Two local calculations govern the comparison. We choose generators
for which sextic reciprocity has no supplementary factor. In a
reflected interaction between a prime occurring to valuation one
in the varying twist and a varying squarefree coefficient prime,
the sextic exponents are $1$ and $-4$. Their sum is $3$ modulo $6$,
so the interaction is quadratic. The quadratic large sieve of
Heath-Brown, in its number-field form due to Goldmakher and Louvel,
then controls the reflected moment
\cite{HeathBrown95,GoldmakherLouvel13}.

On the Poisson side, a relative-norm formula identifies the cube of
the cubic Gauss phase with the negative of a unitary Hecke-character
value. A finite-order correction gives the target character $\eta$.
At a prime ideal $\mathfrak p$ outside a fixed excluded set, the
first squarefree theta term at the principal frequency is
$-\eta(\mathfrak p)(\Norm\mathfrak p)^{-\xi}$, where $\xi$ is
the Mellin variable for the coefficient norm. Together with the
constant term, it supplies the local factor of the reciprocal
target $L$-function. The remaining local terms give a holomorphic,
nonvanishing correction in the region used at the end of the proof.

Reflection and Poisson summation thus estimate the same additive
average in two ways. A conductor large sieve and zero detection
control the nonprincipal frequencies. The resulting power saving
for the principal inverse Mellin integral continues its Mellin
transform across the line $\re s=1$ and excludes zeros of the
target function. The use of ideal norms keeps the numerical
exponents in this argument independent of the degree of $F$.
Sections~\ref{sec:arithmetic}--\ref{sec:zero-free} carry out this
proof; Appendix~\ref{app:whittaker} provides the theta normalization.

\subsection{Controlled predecessors in a fixed progression}

The construction produces many primes of the form
\[
 p=crQ+1,\qquad c\in\{2,4\},\qquad Q>x^{0.9}\text{ prime},
\]
where every prime divisor of $r$ lies between
$\exp((\log x)^{0.1})$ and $\exp((\log x)^{0.3})$.
We also prescribe a residue class on which the Legendre symbol
$(a/p)$ equals $-1$. This excludes $2$ as an index divisor.
A nontrivial remaining index has a prime divisor in $rQ$:
the factor $Q$ leads to a small elementary exceptional set, and
the factors of $r$ fall within the uniform splitting estimate.

The construction starts with weighted integers $d=crQ+1$ of this
shape. A sieve removes small prime factors of $d$. For a remaining
composite $d=mn$, with $n$ its least prime factor, a marked bilinear
estimate replaces the restriction on the cofactor $m$ by an explicit
density and a weaker condition on its small prime factors.
Buchstab's least-prime-factor identity leaves a positive margin
between the initial sifted mass and the mass removed in this
comparison. A separate upper bound makes the contribution from two
prime factors close to $\sqrt x$ smaller than that margin.

The marked prime estimate is an input from
\emph{The Poisson--Dirichlet law for the prime factors of $p-1$}
\cite[Theorem~3.1]{PoissonDirichlet}. We extract its scalar coefficient
criterion and verify the full character and frequency estimates for
the rough-factor comparison. The weights require $d-1$ to have a
prime divisor in each of several prescribed ranges; these choices
are the marks. A bounded multiplier in the bilinear estimate may
depend on the quotient left after every power of these primes has
been removed. A congruence condition on $d=mn$ need not have this
invariance. We therefore expand $mn\equiv u\pmod M$ in Dirichlet
characters and separate it into factors $\lambda(m)\lambda(n)$
in the two scalar coefficient sequences. Fixed twists preserve
the required long character-sum bound.

The resulting fixed-progression construction is stated separately
in Proposition~\ref{prop:controlled}. It applies to any fixed data
$M,c,u$ satisfying its local coprimality conditions, independently
of the primitive-root application. Its proof uses
Bombieri--Vinogradov distribution, a Brun--Hooley block sieve,
and the continuous rough-number density
\cite{Bombieri65,Vinogradov65,FordHalberstam2000,Buchstab1937}.
The precise elementary sieve and density forms are adapted from
\cite[Lemmas~2.9 and~7.4]{BlockPaper}, with proofs included here.

\subsection{Organization and notation}

Section~\ref{sec:reduction} states the controlled-predecessor and
splitting estimates and combines them to prove
Theorem~\ref{thm:primitive-root}. Sections~\ref{sec:arithmetic}--\ref{sec:zero-free}
prove Theorem~\ref{thm:hecke-zero-free} for a fixed arbitrary $F$
and $\eta$, independently of the integer base. The analytic scale
parameters in those sections are local to that proof.
Section~\ref{sec:splitting} returns to the fixed integer $a$ and
derives the uniform splitting bound. Section~\ref{sec:type-ii}
proves the marked coefficient and rough-factor estimates;
Section~\ref{sec:sieve} supplies the elementary sieve facts;
Section~\ref{sec:prime-construction} proves the fixed-progression
construction. Appendix~\ref{app:whittaker} verifies the selected
Kazhdan--Patterson specialization and Whittaker normalization.

Throughout, $\log$ is the natural logarithm, and $L=\log x$ in the
prime-counting arguments. The symbols $\ll$ and
$O$ have their usual meaning; subscripts record permitted dependence
when it matters. Auxiliary parameters are fixed before the relevant
asymptotic variable tends to infinity. Uniformity in those parameters
is stated explicitly. The general representation-theoretic,
class-field, Hecke-continuation, and classical sieve theorems are
cited at their points of use. The marked Type II argument also uses
the companion papers cited in the bibliography.
